\section{Results}
\label{sec:results}

\subsection{Residual Stream Similarity: Compounds Are Modifications, Not New Directions}
\label{sec:results_exp1}

\tabref{tab:cosine} presents the cosine similarity between compound and component representations at the final layer (layer 11).
The mean similarity between compound and head noun representations is $0.954 \pm 0.019$, and between compound and additive composition is $0.959 \pm 0.013$.
Even the least compositional compound in our dataset, ``red herring,'' maintains a cosine similarity of 0.895 with its head noun alone.

\begin{table}[t]
    \centering
    \resizebox{\textwidth}{!}{%
    \begin{tabular}{@{}lccccc@{}}
        \toprule
        \textbf{Compound} & \textbf{Comp.\ Level} & \textbf{vs.\ Head Alone} & \textbf{vs.\ Additive} & \textbf{vs.\ Modifier} \\
        \midrule
        school bus & high & {\bf 0.983} & 0.966 & 0.823 \\
        swimming pool & high & 0.961 & 0.965 & 0.810 \\
        coffee machine & high & 0.957 & 0.967 & 0.844 \\
        kitchen chair & very\_high & 0.956 & 0.968 & 0.831 \\
        washing machine & medium & 0.951 & 0.954 & 0.812 \\
        guinea pig & low & 0.942 & 0.938 & 0.760 \\
        ice cream & medium & 0.937 & 0.943 & 0.762 \\
        hot dog & low & 0.936 & 0.943 & 0.761 \\
        dark horse & very\_low & 0.936 & 0.947 & 0.777 \\
        red herring & very\_low & 0.895 & 0.925 & 0.748 \\
        \midrule
        \textbf{Mean (all 21)} & --- & $0.954$ {\scriptsize $\pm 0.019$} & $0.959$ {\scriptsize $\pm 0.013$} & --- \\
        \bottomrule
    \end{tabular}
    }
    \caption{Cosine similarity between compound representations and component representations at layer 11. All compounds maintain high similarity ($>0.89$) with their head nouns, indicating that compound representations are modifications of existing component directions rather than new directions. Best per-row result in \textbf{bold} when noteworthy.}
    \label{tab:cosine}
\end{table}

{\bf Compositionality predicts residual stream similarity.}
There is a significant positive correlation between compositionality and similarity with the additive composition baseline (Spearman $\rho = 0.602$, $p = 0.004$): transparent compounds like ``kitchen chair'' (0.968) are closer to the additive baseline than idiomatic compounds like ``red herring'' (0.925).
However, even the most idiomatic compounds show very high cosine similarity, suggesting that residual stream geometry provides a limited window into compound-specific processing.

\subsection{SAE Features Reveal Hidden Compound-Specific Computation}
\label{sec:results_exp2}

Despite high cosine similarity at the residual stream level, SAE decomposition reveals that a majority of features active for compounds are not shared with either component alone.
\tabref{tab:sae} presents the SAE feature analysis at layer 11.
The mean fraction of unique-to-compound features is $0.563 \pm 0.131$, and the mean Jaccard similarity between compound and head feature sets is only $0.255 \pm 0.091$.

\begin{table}[t]
    \centering
    \resizebox{\textwidth}{!}{%
    \begin{tabular}{@{}lcccccc@{}}
        \toprule
        \textbf{Compound} & \textbf{Comp.} & \textbf{Unique (\%)} & \textbf{Head Overlap (\%)} & \textbf{Mod.\ Overlap (\%)} & \textbf{Jaccard w/ Head} \\
        \midrule
        red herring & very\_low & {\bf 82.1} & 14.3 & 7.1 & 0.079 \\
        dark horse & very\_low & 78.3 & 19.3 & 13.3 & 0.147 \\
        hot dog & low & 66.7 & 23.5 & 15.7 & 0.152 \\
        guinea pig & low & 61.7 & 31.7 & 18.3 & 0.209 \\
        washing machine & medium & 57.9 & 35.1 & 24.6 & 0.256 \\
        coffee machine & high & 54.7 & 34.4 & 20.3 & 0.265 \\
        swimming pool & high & 53.7 & 33.3 & 20.4 & 0.231 \\
        school bus & high & 48.9 & 46.7 & 28.9 & 0.333 \\
        kitchen chair & very\_high & 46.9 & 34.7 & 32.7 & 0.205 \\
        dish washer & high & 21.5 & {\bf 69.2} & 27.7 & {\bf 0.479} \\
        \midrule
        \textbf{Mean (all 21)} & --- & $56.3$ {\scriptsize $\pm 13.1$} & --- & --- & $0.255$ {\scriptsize $\pm 0.091$} \\
        \bottomrule
    \end{tabular}
    }
    \caption{SAE feature decomposition at layer 11. ``Unique'' is the fraction of compound features not shared with either component alone. ``Head/Mod.\ Overlap'' is the fraction of compound features also active for the head/modifier in isolation. Idiomatic compounds activate substantially more unique features than transparent ones. Best (most extreme) values in \textbf{bold}.}
    \label{tab:sae}
\end{table}

{\bf Compositionality strongly predicts unique feature fraction.}
Low-compositionality compounds (very\_low and low) have significantly more unique SAE features (mean 71.2\%) than high-compositionality compounds (mean 48.2\%).
A Mann-Whitney U test confirms this difference is statistically significant ($p = 0.0005$), with a very large effect size (Cohen's $d = 2.54$).
This makes linguistic sense: idioms like ``red herring'' require more specialized representation because their meaning cannot be derived from their components.

{\bf ``Dish washer'' is an informative outlier.}
With only 21.5\% unique features and 69.2\% overlap with ``washer'' alone, ``dish washer'' is the most compositional compound by SAE metrics---the model treats it almost identically to ``washer.''
This is semantically reasonable: a dishwasher \emph{is} a type of washer, so minimal additional representation is needed.

{\bf Cosine similarity versus SAE decomposition.}
The contrast between these two analysis levels is striking.
``Red herring'' has 0.895 cosine similarity with ``herring'' alone (seemingly very similar), yet only 14.3\% feature overlap.
This demonstrates that cosine similarity measures overall direction alignment while SAE decomposition captures fine-grained feature-level differences invisible to coarser geometric measures.

\subsection{Next-Token Priming: The Primary Storage Mechanism}
\label{sec:results_exp3}

\tabref{tab:priming} presents the next-token prediction priming results.
The variation across compounds is enormous, spanning five orders of magnitude from $1.3\times$ (``office desk'') to $28{,}037\times$ (``guinea pig'').

\begin{table}[t]
    \centering
    \resizebox{\textwidth}{!}{%
    \begin{tabular}{@{}lccccc@{}}
        \toprule
        \textbf{Compound} & $P(\text{head} \mid \text{mod})$ & $P(\text{head} \mid \text{baseline})$ & \textbf{Priming Ratio} & \textbf{Rank} \\
        \midrule
        guinea pig & 0.6624 & 0.000024 & ${\bf 28{,}037}\times$ & {\bf 1} \\
        vending machine & 0.5416 & 0.000112 & $4{,}853\times$ & {\bf 1} \\
        washing machine & 0.4711 & 0.000112 & $4{,}221\times$ & {\bf 1} \\
        swimming pool & 0.3184 & 0.000079 & $4{,}022\times$ & {\bf 1} \\
        sewing machine & 0.3741 & 0.000112 & $3{,}352\times$ & {\bf 1} \\
        ice cream & 0.0743 & 0.000020 & $3{,}757\times$ & 3 \\
        slot machine & 0.0624 & 0.000112 & $559\times$ & 2 \\
        hot dog & 0.0307 & 0.000127 & $242\times$ & 5 \\
        dark horse & 0.0097 & 0.000067 & $144\times$ & 8 \\
        kitchen chair & 0.0012 & 0.000064 & $19\times$ & 83 \\
        red herring & 0.0019 & 0.000434 & $4.3\times$ & 67 \\
        rain coat & 0.0001 & 0.000015 & $3.9\times$ & 888 \\
        office desk & 0.0000 & 0.000017 & $1.3\times$ & 1429 \\
        \bottomrule
    \end{tabular}
    }
    \caption{Next-token prediction priming results. The priming ratio (\eqref{eq:priming}) measures how much more likely the head noun is after the modifier compared to baseline. ``Frozen'' compounds like ``guinea pig'' and ``washing machine'' show extreme priming ($>4{,}000\times$), while loosely associated compounds like ``office desk'' show minimal priming ($1.3\times$). Rank indicates the head noun's position among all 50,257 vocabulary predictions.}
    \label{tab:priming}
\end{table}

{\bf Three regimes of compound priming.}
The data naturally cluster into three groups:
\begin{enumerate}[leftmargin=*,itemsep=0pt,topsep=0pt]
    \item \emph{Extreme priming} ($>1{,}000\times$): washing machine, vending machine, sewing machine, swimming pool, guinea pig. These are ``frozen'' compounds where the modifier almost exclusively precedes its head in natural text. The head noun is the \#1 prediction.
    \item \emph{Moderate priming} ($10$--$1{,}000\times$): hot dog, slot machine, dark horse, ice cream. The modifier has other common continuations, but the compound is still strongly predicted.
    \item \emph{Minimal priming} ($<10\times$): red herring ($4.3\times$), rain coat ($3.9\times$), office desk ($1.3\times$). These modifiers have many common continuations and the compound is either rare or not strongly associated.
\end{enumerate}

{\bf ``Washing machine'' is constructed through prediction.}
After ``The washing,'' GPT-2 produces the following top-5 predictions: \texttt{machine} (47.1\%), \texttt{of} (8.4\%), \texttt{-} (8.2\%), \texttt{machines} (4.9\%), \texttt{up} (3.8\%).
The model assigns nearly half its probability mass to ``machine,'' effectively constructing the compound concept through its transition probabilities.
This is the primary mechanism by which GPT-2 ``stores'' the concept of a washing machine: not as a dedicated direction, but as an overwhelming statistical association.

{\bf ``Guinea pig'' is the most frozen compound.}
Despite being rated as low compositionality (the meaning of ``guinea pig'' is opaque), ``guinea'' predicts ``pig'' with 66.2\% probability and a $28{,}037\times$ priming ratio.
This dissociation between semantic compositionality and statistical predictability reveals that priming reflects co-occurrence frequency, not semantic transparency.
